\cgeentry{cge_dynresults}{Results of a dynamic run as a dataset}

\cgesec{Description}

\cmd{cge_dynresults} replaces the data in memory with the results of a run
solved by \manref{cge_dynamic}: one observation for each period and each
element of every model variable, with its base value and its value in the
period. With \cmd{baseline()}, the value of a baseline run in the same period
and the per cent difference from it are added. The dataset can then be
listed, tabulated, graphed or saved with any Stata command.

\cgesec{Quick start}

Load every variable of every period of run \cmd{base}
\begin{stcode}
cge_dynresults, run(base) clear
\end{stcode}
Load the periods of run \cmd{exit} with their differences from run \cmd{base},
and graph the per cent difference in the consumption of household
\cmd{hh_ru_p}
\begin{stcode}
cge_dynresults, run(exit) baseline(base) clear
line dev_pct period if variable == "EH" & element1 == "hh_ru_p"
\end{stcode}

\cgesec{Syntax}

\begin{stsyntax}
\cmdb{cge_dynresults}\texttt{,} \optarg{}{run}{name} \opt{}{clear}
[\optarg{base}{line}{name}]
\end{stsyntax}

\begin{opttable}
\optgroup{Main}
\optarg{}{run}{name} & the dynamic run; required\\
\opt{}{clear} & replace the data in memory; required\\
\optarg{base}{line}{name} & the baseline run to compare with\\
\end{opttable}

\cgesec{Options}

\begin{optdesc}
\item[\optarg{}{run}{name}] names the run, as given in \cmd{name()} of
\manref{cge_dynamic}.

\item[\opt{}{clear}] is required: the data in memory are replaced. They are
kept if the run cannot be loaded, for example because no run has the name
given.

\item[\optarg{base}{line}{name}] names a baseline run to compare with. It
must have every period of \cmd{run()}; it may have more.
\end{optdesc}

\cgesec{Remarks and examples}

The dataset has one observation for each period of the run and each element
of every model variable that exists in the model, in period order. A variable
is identified by its name, as listed in \secref{sec:variables}, and by up to
two set elements: \cmd{QA} has one, the activity; \cmd{QF} has two, the
factor and the activity; \cmd{EXR} has none. The variables of the dataset are

\begin{cgetable}{>{\ttfamily}l p{0.78\linewidth}}
\toprule
\normalfont Variable & Content\\
\midrule\endfirsthead
\toprule
\normalfont Variable & Content\\
\midrule\endhead
\bottomrule\endlastfoot
period & the period\\
variable & the model variable\\
element1 & its first set element, or empty\\
element2 & its second set element, or empty\\
base & its base value, in the calibrated SAM\\
value & its value in the period\\
baseline & with \cmd{baseline()}: its value in the baseline run in the same period\\
dev\_pct & with \cmd{baseline()}: $100\,(\text{value} - \text{baseline})/\lvert\text{baseline}\rvert$, missing where the baseline value is zero\\
\end{cgetable}

The data are sorted by \cmd{period}, \cmd{variable}, \cmd{element1} and
\cmd{element2}. \manref{cge_report} with \cmd{run()} gives the standard report
measures of each period as changes from the base; \cmd{cge_dynresults} gives
every variable, and its difference from a baseline in each period.

\cgesubsec{Example 1: a scenario against its baseline}

The two runs of Examples 1 and 2 of \manref{cge_dynamic} are solved first, on
the synthetic SAM: a baseline \cmd{base} calibrated to 4 per cent growth of
real value added, and a scenario \cmd{lib} that removes import tariffs from
period 2, with the baseline's efficiency path:
\begin{stcode}
. cge_closure, factors(cap:specific)
. local settings capital(cap) depreciation(0.05) rental(0.15) growth(lab:0.025) ///
>     population(0.025) govgrowth(0.03) foreigngrowth(0.02)
. cge_dynamic, name(base) periods(0/4) `settings' gdpgrowth(0.04)
. cge_dynamic, name(lib) periods(0/4) `settings' baseline(base) shocks(tariff_from2)
\end{stcode}
The results of the scenario, with the baseline's, replace the data in memory:
\stlog{dynr_load}
Capital is specific to each activity. Its stocks are the same in both runs up
to period 2, the period of the shock, because capital is set before each
period from the investment of the period before; from period 3 they are lower
in the scenario:
\stlog{dynr_capital}
Output changes in the period of the shock; output of services rises above the
baseline while that of agriculture and manufactures falls:
\stlog{dynr_output}
The SAM is synthetic: these results illustrate the command and are not
findings about any economy.

\cgesubsec{Example 2: a run on its own}

Without \cmd{baseline()}, the dataset has the base value and the value in
each period:
\stlog{dynr_nobase}

\cgesec{Stored results}

\cmd{cge_dynresults} stores the following in \cmd{r()}:
\begin{storedres}
\resgroup{Macros}
\res{r(run)} & name of the run\\
\end{storedres}

\cgesec{Methods and formulas}

For element $i$ of a variable in period $t$, with value $x_{it}$ in the run
and $b_{it}$ in the baseline run,
\begin{equation*}
\text{\cmd{dev_pct}}_{it} = 100\,\frac{x_{it} - b_{it}}{\lvert b_{it}\rvert}, \qquad b_{it} \ne 0,
\end{equation*}
and missing when $b_{it} = 0$. The base value is the value of the variable
in the calibrated SAM, the same in every period. \Chapref{chap:dynamics}
describes the runs.

\cgesec{Also see}

\begin{alsosee}
\seeentry{cge_dynamic}{Solve a recursive dynamic run}
\seeentry{cge_report}{Standard report tables for one or more simulations}
\seeentry{cge_results}{Report simulation results}
\seechapter{chap:dynamics}{Recursive dynamics}
\end{alsosee}
